%!TeX root=MemoriaTFG.tex

\chapter{Summary}

The purpose of this report is to present the Wave Launcher; a local system for managing Minecraft servers, and compare it against the traditional way of managing said servers.
There are few Minecraft server launchers that are both simple enough for inexperienced users and capable of meeting their most common needs. The Wave Launcher attempts to fill this niche by providing a simple user interface that anyone can understand while internally managing the server’s state to ensure that it remains correctly configured. It is designed to manage Minecraft: Java Edition servers and currently targets Windows; however, its architecture and implementation language (C\#) were chosen to facilitate porting it to other platforms, including macOS, Linux, and potentially a web interface.
For the user interface, the application uses .NET MAUI 10, as it provides cross-platform support while rendering native controls on each supported platform. The application follows an architecture inspired by the hexagonal architecture, separating responsibilities to keep the codebase as clean and maintainable as possible. Unlike a conventional web-based implementation, this adapted architecture does not include a separate API layer because the UI and application services are both written in C\# and run within the same application on the same machine.
Managing a Minecraft server requires coordinating several interdependent components, including the Java installation used to run the server, the mod loader responsible for loading modifications, and the mods themselves. The Wave Launcher connects to the APIs exposed by the selected providers to retrieve and manage these components. It supports the Java runtime distributed by Mojang and Eclipse Adoptium, the Forge and Fabric mod loaders, and mods obtained from Modrinth and CurseForge.